\documentclass{article}
\usepackage{amsmath,amssymb,amsthm}
\usepackage{algorithm}
\usepackage[noend]{algpseudocode}
\usepackage{booktabs}
\usepackage{enumitem}
\usepackage{hyperref}
\usepackage{xcolor}
\newcommand{\E}{\mathbb{E}}
\newcommand{\R}{\mathbb{R}}
\newtheorem{theorem}{Theorem}
\newtheorem{lemma}{Lemma}
\newtheorem{assumption}{Assumption}
\newtheorem{corollary}{Corollary}
\newtheorem{definition}{Definition}
\begin{document}

\title{Local SGD with Proximal Regularization (FedProx) \\ Detailed Algorithmic Description and Convergence Proof}
\author{}
\date{}
\maketitle

%%%%%%%%%%%%%%%%%%%%%%%%%%%%%%%%%%%%%%%%%%%%%%%%%%%%%%%%%%%%%%%%%%%%%%%%%%%%%%
\section*{Algorithm Name}
\textbf{FedProx: Local Stochastic Gradient Descent with a Proximal Term}  

The method performs $E$ local SGD steps on each of $K$ participating clients.  
Before each local step the client solves a proximal‐regularized subproblem that
penalizes deviation from the current global model $x^{(t)}$.

%%%%%%%%%%%%%%%%%%%%%%%%%%%%%%%%%%%%%%%%%%%%%%%%%%%%%%%%%%%%%%%%%%%%%%%%%%%%%%
\section*{Mathematical Formulation}
\noindent
\textbf{Global objective.}  
Each client $k\in\{1,\dots,K\}$ possesses a smooth (possibly non‑convex) loss
\[
f_k(x)\;=\;\frac{1}{n_k}\sum_{i=1}^{n_k} \ell\!\bigl(x;z_{k,i}\bigr),\qquad 
x\in\R^d .
\]
The federated objective is the weighted average
\[
F(x)=\sum_{k=1}^{K}p_k f_k(x),\qquad p_k:=\frac{n_k}{\sum_{j=1}^{K}n_j},\;
\sum_k p_k =1 .
\]

\medskip
\textbf{Local proximal subproblem.}  
At communication round $t$ the server broadcasts the global model $x^{(t)}$.
Client $k$ solves (approximately) for $e=0,\dots,E-1$
\begin{equation}
\label{eq:local_update}
x_{k}^{(t,e+1)} \;=\;
x_{k}^{(t,e)} \;-\;\eta_{t}\Bigl(\nabla \ell\!\bigl(x_{k}^{(t,e)}; \xi_{k}^{(t,e)}\bigr)
\;+\;\mu\bigl(x_{k}^{(t,e)}-x^{(t)}\bigr)\Bigr),
\end{equation}
where
\begin{itemize}
    \item $\eta_{t}>0$ is the (possibly diminishing) learning rate,
    \item $\xi_{k}^{(t,e)}$ is a stochastic sample drawn uniformly from client $k$’s data,
    \item $\mu\ge0$ is the proximal coefficient.
\end{itemize}
The initial local point is $x_{k}^{(t,0)}=x^{(t)}$.

\medskip
\textbf{Aggregation.}  After $E$ local steps each client returns $x_{k}^{(t,E)}$.
The server updates the global model by weighted averaging:
\begin{equation}
\label{eq:global_agg}
x^{(t+1)}\;=\;\sum_{k=1}^{K} p_k\,x_{k}^{(t,E)} .
\end{equation}

\medskip
\textbf{Hyper‑parameters.}
\begin{itemize}
    \item $K$: number of selected clients per round (we assume all $K$ are used).
    \item $E$: number of local SGD steps between two aggregations.
    \item $\eta_{t}$: learning‑rate schedule (e.g. $\eta_{t}= \eta_0/\sqrt{t+1}$).
    \item $\mu$: proximal penalty weight (FedProx parameter).
\end{itemize}

%%%%%%%%%%%%%%%%%%%%%%%%%%%%%%%%%%%%%%%%%%%%%%%%%%%%%%%%%%%%%%%%%%%%%%%%%%%%%%
\section*{Pseudocode}
\begin{algorithm}[H]
\caption{FedProx (Local SGD with Proximal Term)}
\begin{algorithmic}[1]
\State \textbf{Input:} initial model $x^{(0)}$, learning‑rate schedule $\{\eta_t\}$,
proximal coefficient $\mu\ge0$, local steps $E$, total rounds $T$.
\For{$t=0,1,\dots,T-1$}
    \State Server broadcasts $x^{(t)}$ to all clients.
    \For{each client $k=1,\dots,K$ \textbf{in parallel}}
        \State $x_{k}^{(t,0)} \gets x^{(t)}$
        \For{$e=0$ \textbf{to} $E-1$}
            \State Sample a minibatch $\xi_{k}^{(t,e)}$ from local data.
            \State Compute stochastic gradient $g_{k}^{(t,e)}=
                    \nabla \ell\!\bigl(x_{k}^{(t,e)};\xi_{k}^{(t,e)}\bigr)$.
            \State \textbf{Local update:}
            \[
            x_{k}^{(t,e+1)} \gets
            x_{k}^{(t,e)} - \eta_{t}\bigl(g_{k}^{(t,e)}+\mu(x_{k}^{(t,e)}-x^{(t)})\bigr).
            \]
        \EndFor
        \State Send $x_{k}^{(t,E)}$ to server.
    \EndFor
    \State Server aggregates:
    \[
    x^{(t+1)} \gets \sum_{k=1}^{K} p_k\,x_{k}^{(t,E)} .
    \]
\EndFor
\State \textbf{Output:} final model $x^{(T)}$.
\end{algorithmic}
\end{algorithm}

%%%%%%%%%%%%%%%%%%%%%%%%%%%%%%%%%%%%%%%%%%%%%%%%%%%%%%%%%%%%%%%%%%%%%%%%%%%%%%
\section*{Dry Run Example}
Consider a single‑client setting ($K=1$) with scalar parameter $w\in\R$ and loss
\[
f(w)= (w-3)^2 .
\]
Thus $\nabla f(w)=2(w-3)$.  
Take $w^{(0)}=0$, learning rate $\eta=0.1$, proximal coefficient $\mu=0.5$,
and $E=3$ local steps (no communication needed because $K=1$).

\medskip
\begin{tabular}{@{}c|c|c|c@{}}
\toprule
Step & $w_{k}^{(e)}$ & Gradient $g^{(e)}$ & Updated $w_{k}^{(e+1)}$ \\
\midrule
0 & $0$ & $2(0-3)=-6$ & $0-0.1\bigl(-6+0.5(0-0)\bigr)=0.6$ \\
1 & $0.6$ & $2(0.6-3)=-4.8$ & $0.6-0.1\bigl(-4.8+0.5(0.6-0)\bigr)=0.6-0.1(-4.8+0.3)=0.6+0.45=1.05$ \\
2 & $1.05$ & $2(1.05-3)=-3.9$ & $1.05-0.1\bigl(-3.9+0.5(1.05-0)\bigr)=1.05-0.1(-3.9+0.525)=1.05+0.3375=1.3875$ \\
\bottomrule
\end{tabular}

Objective values after each local step:
\[
f(0)=9,\; f(0.6)=5.76,\; f(1.05)=3.8025,\; f(1.3875)=2.5955.
\]
The sequence shows descent toward the minimizer $w^\star=3$.

%%%%%%%%%%%%%%%%%%%%%%%%%%%%%%%%%%%%%%%%%%%%%%%%%%%%%%%%%%%%%%%%%%%%%%%%%%%%%%
\section*{Assumptions}
\begin{assumption}[Smoothness]
Each local loss $\ell(x;z)$ is $L$‑smooth, i.e.
\[
\|\nabla \ell(x;z)-\nabla \ell(y;z)\|\le L\|x-y\|,\qquad\forall x,y\in\R^d,\;z.
\]
Consequently each $f_k$ is $L$‑smooth.
\end{assumption}

\begin{assumption}[Bounded Stochastic Gradient Variance]
For every client $k$ and any $x$,
\[
\E_{\xi\sim\mathcal{D}_k}\!\bigl[\|\nabla \ell(x;\xi)-\nabla f_k(x)\|^2\bigr]\le\sigma^2 .
\]
\end{assumption}

\begin{assumption}[Lower Bounded Objective]
The global objective $F$ is bounded below: $F^\star:=\inf_{x}F(x)>-\infty$.
\end{assumption}

\begin{assumption}[Learning Rate]
The stepsize satisfies
\[
\eta_t \le \frac{1}{L+\mu}\quad\text{for all }t,
\]
and $\sum_{t=0}^{\infty}\eta_t = \infty$, $\sum_{t=0}^{\infty}\eta_t^{2}<\infty$ 
(e.g. $\eta_t=\eta_0/(t+1)^{\alpha}$ with $\alpha\in(0.5,1]$).
\end{assumption}

\begin{assumption}[Client Participation]
All $K$ clients participate in every round (full participation).  
Thus $p_k=n_k/\sum_j n_j$ are fixed probabilities.
\end{assumption}

%%%%%%%%%%%%%%%%%%%%%%%%%%%%%%%%%%%%%%%%%%%%%%%%%%%%%%%%%%%%%%%%%%%%%%%%%%%%%%
\section*{Main Convergence Theorem}
\begin{theorem}[Convergence to Stationarity]\label{thm:main}
Under Assumptions 1–5, let $\{x^{(t)}\}_{t\ge0}$ be the sequence generated by
FedProx with $E\ge1$ local steps. Define the averaged squared gradient norm
\[
\Phi_T\;:=\;\frac{1}{T}\sum_{t=0}^{T-1}\E\bigl[\|\nabla F(x^{(t)})\|^2\bigr].
\]
Then for any $T\ge1$,
\[
\Phi_T\;\le\;
\frac{2\bigl(F(x^{(0)})-F^\star\bigr)}{\eta_{\min}T}
\;+\;
\frac{2L\sigma^2\eta_{\max}E}{K}
\;+\;
\frac{2\mu^2 L\eta_{\max}E^2}{K},
\]
where $\eta_{\min}:=\min_{t<T}\eta_t$ and $\eta_{\max}:=\max_{t<T}\eta_t$.
Consequently, if $\eta_t=\Theta(1/\sqrt{T})$, we obtain
\[
\Phi_T = O\!\left(\frac{1}{\sqrt{T}}\right).
\]
\end{theorem}

\begin{corollary}[Special case $\mu=0$ (plain Local SGD)]
When $\mu=0$ the bound reduces to the standard Local SGD rate
\[
\Phi_T\le
\frac{2\bigl(F(x^{(0)})-F^\star\bigr)}{\eta_{\min}T}
+\frac{2L\sigma^2\eta_{\max}E}{K}.
\]
\end{corollary}

%%%%%%%%%%%%%%%%%%%%%%%%%%%%%%%%%%%%%%%%%%%%%%%%%%%%%%%%%%%%%%%%%%%%%%%%%%%%%%
\section*{Theoretical Properties}
\begin{itemize}
\item \textbf{Stability w.r.t. Heterogeneity.} The proximal term $\mu\|x-x^{(t)}\|^2$ controls client drift; larger $\mu$ yields a tighter bound on the drift term $ \mu^2 L\eta_{\max}E^2/K$ at the expense of slower local progress.
\item \textbf{Hyper‑parameter Sensitivity.}
    \begin{enumerate}[label=(\alph*)]
        \item \emph{Learning rate $\eta_t$}: must respect $\eta_t\le 1/(L+\mu)$ for the descent Lemma; diminishing schedules guarantee the summability conditions.
        \item \emph{Local steps $E$}: appears linearly in the variance term and quadratically in the drift term; moderate $E$ balances communication efficiency and convergence speed.
        \item \emph{Proximal coefficient $\mu$}: when data are highly non‑i.i.d., taking $\mu>0$ reduces the $E^2$ drift penalty.
    \end{enumerate}
\item \textbf{Comparison to Baselines.}
    \begin{itemize}
        \item With $\mu=0$, Theorem~\ref{thm:main} recovers the classical $O(1/\sqrt{T})$ rate for non‑convex Local SGD.
        \item For heterogeneous data, the additional $O(\mu^2E^2)$ term is usually smaller than the $O(E^2)$ drift that appears in naive Local SGD analyses, explaining the empirical robustness of FedProx.
    \end{itemize}
\end{itemize}

%%%%%%%%%%%%%%%%%%%%%%%%%%%%%%%%%%%%%%%%%%%%%%%%%%%%%%%%%%%%%%%%%%%%%%%%%%%%%%
\section*{Preliminaries (Definitions and Lemmas)}
\begin{definition}[Full‑gradient and Local Gradient]
For any $x$, define the global gradient $\nabla F(x)=\sum_{k}p_k\nabla f_k(x)$.
For client $k$ at local iteration $(t,e)$, the stochastic gradient is
$g_{k}^{(t,e)}:=\nabla\ell(x_{k}^{(t,e)};\xi_{k}^{(t,e)})$.
\end{definition}

\begin{lemma}[Smoothness Descent Lemma]\label{lem:smooth}
If $h$ is $L$‑smooth, then for any $u,v$,
\[
h(v)\le h(u)+\langle\nabla h(u),v-u\rangle+\frac{L}{2}\|v-u\|^2 .
\]
\end{lemma}

\begin{lemma}[Unbiased Gradient and Bounded Variance]\label{lem:unbiased}
For any $k$ and $x$, $\E[g_{k}^{(t,e)}\mid x_{k}^{(t,e)}]=\nabla f_k(x_{k}^{(t,e)})$ and
$\E\|g_{k}^{(t,e)}-\nabla f_k(x_{k}^{(t,e)})\|^2\le\sigma^2$.
\end{lemma}

%%%%%%%%%%%%%%%%%%%%%%%%%%%%%%%%%%%%%%%%%%%%%%%%%%%%%%%%%%%%%%%%%%%%%%%%%%%%%%
\section*{Main Proof (Step‑by‑Step Derivation)}
We prove Theorem~\ref{thm:main}. Throughout the proof we condition on the history
$\mathcal{F}_t:=\sigma\bigl(x^{(0)},\{g_{k}^{(s,e)}\}_{s<t}\bigr)$.

\medskip
\noindent\textbf{[Step 1] Descent of the global objective after one communication round.}  
Using Lemma~\ref{lem:smooth} with $h=F$, $u=x^{(t)}$, $v=x^{(t+1)}$ and the fact that
$x^{(t+1)}=\sum_k p_k x_{k}^{(t,E)}$, we have
\begin{align}
F(x^{(t+1)}) 
&\le F(x^{(t)}) 
   +\bigl\langle\nabla F(x^{(t)}),\,x^{(t+1)}-x^{(t)}\bigr\rangle
   +\frac{L}{2}\|x^{(t+1)}-x^{(t)}\|^2 . \tag{1}
\end{align}

\medskip
\noindent\textbf{[Step 2] Expand $x^{(t+1)}-x^{(t)}$ using the local updates.}  
From (\ref{eq:local_update}) and telescoping over $e$,
\[
x_{k}^{(t,E)}-x^{(t)} 
 = -\eta_t\sum_{e=0}^{E-1}\bigl(g_{k}^{(t,e)}+\mu(x_{k}^{(t,e)}-x^{(t)})\bigr).
\]
Hence
\[
x^{(t+1)}-x^{(t)} = -\eta_t\sum_{k=1}^{K}p_k
      \sum_{e=0}^{E-1}\bigl(g_{k}^{(t,e)}+\mu(x_{k}^{(t,e)}-x^{(t)})\bigr). \tag{2}
\]

\medskip
\noindent\textbf{[Step 3] Plug (2) into (1) and take expectation conditioned on $\mathcal{F}_t$.}  
Using linearity of expectation and Lemma~\ref{lem:unbiased},
\begin{align}
\E\!\bigl[F(x^{(t+1)})\mid\mathcal{F}_t\bigr]
&\le F(x^{(t)}) 
   -\eta_t\underbrace{
      \Bigl\langle\nabla F(x^{(t)}),\,
      \sum_{k}p_k\sum_{e=0}^{E-1}
      \bigl(\nabla f_k(x_{k}^{(t,e)})+\mu(x_{k}^{(t,e)}-x^{(t)})\bigr)
      \Bigr\rangle}_{\text{Term A}}
   \nonumber\\
&\quad +\frac{L\eta_t^{2}}{2}\underbrace{
      \E\!\Bigl[\bigl\|
      \sum_{k}p_k\sum_{e=0}^{E-1}
      (g_{k}^{(t,e)}+\mu(x_{k}^{(t,e)}-x^{(t)}))
      \bigr\|^{2}
      \,\big|\,\mathcal{F}_t\Bigr]}_{\text{Term B}} .
\tag{3}
\end{align}

\medskip
\noindent\textbf{[Step 4] Lower‑bound Term A.}  
Write $\nabla F(x^{(t)})=\sum_k p_k \nabla f_k(x^{(t)})$.
Then
\begin{align}
\text{Term A}
&=\sum_{k}p_k\sum_{e=0}^{E-1}
   \bigl\langle\nabla F(x^{(t)}),\,
          \nabla f_k(x_{k}^{(t,e)})+\mu(x_{k}^{(t,e)}-x^{(t)})\bigr\rangle \nonumber\\
&=\sum_{k}p_k\sum_{e=0}^{E-1}
   \Bigl[
        \langle\nabla F(x^{(t)}),\nabla f_k(x_{k}^{(t,e)})\rangle
        +\mu\langle\nabla F(x^{(t)}),x_{k}^{(t,e)}-x^{(t)}\rangle
   \Bigr]. \tag{4}
\end{align}
Using the identity $2\langle a,b\rangle = \|a\|^2+\|b\|^2-\|a-b\|^2$ and the
smoothness of each $f_k$, we obtain (see Lemma 4.1 in \cite{li2020federated})
\[
\langle\nabla F(x^{(t)}),\nabla f_k(x_{k}^{(t,e)})\rangle
\ge\frac{1}{2}\|\nabla F(x^{(t)})\|^2
      -\frac{L^2}{2}\|x_{k}^{(t,e)}-x^{(t)}\|^2 .
\tag{5}
\]
For the proximal part, Cauchy–Schwarz and $\| \nabla F(x^{(t)})\|\le G$ (bounded
gradient, which follows from $L$‑smoothness and lower boundedness) give
\[
\mu\langle\nabla F(x^{(t)}),x_{k}^{(t,e)}-x^{(t)}\rangle
\ge -\frac{\mu}{2}\|\nabla F(x^{(t)})\|^2
      -\frac{\mu}{2}\|x_{k}^{(t,e)}-x^{(t)}\|^2 .
\tag{6}
\]
Plugging (5)–(6) into (4) and summing over $e$, we obtain
\[
\text{Term A}
\ge
\frac{E}{2}\|\nabla F(x^{(t)})\|^2
-\frac{L^2+\mu}{2}\sum_{k}p_k\sum_{e=0}^{E-1}
   \|x_{k}^{(t,e)}-x^{(t)}\|^2 .
\tag{7}
\]

\medskip
\noindent\textbf{[Step 5] Upper‑bound Term B.}  
Expanding the squared norm and using independence of stochastic gradients across
clients,
\begin{align}
\text{Term B}
&=\E\Bigl[\bigl\|\sum_{k}p_k\sum_{e=0}^{E-1}
      (g_{k}^{(t,e)}-\nabla f_k(x_{k}^{(t,e)}) )
      +\sum_{k}p_k\sum_{e=0}^{E-1}
      (\nabla f_k(x_{k}^{(t,e)})+\mu(x_{k}^{(t,e)}-x^{(t)}))
      \bigr\|^{2}\!\Big|\mathcal{F}_t\Bigr] \nonumber\\
&\le 2\E\Bigl[\bigl\|\sum_{k}p_k\sum_{e=0}^{E-1}
      (g_{k}^{(t,e)}-\nabla f_k(x_{k}^{(t,e)}) )\bigr\|^{2}\!\Big|\mathcal{F}_t\Bigr]
   +2\Bigl\|\sum_{k}p_k\sum_{e=0}^{E-1}
      (\nabla f_k(x_{k}^{(t,e)})+\mu(x_{k}^{(t,e)}-x^{(t)}))\Bigr\|^{2}.
\tag{8}
\end{align}
For the variance part, Jensen’s inequality and Lemma~\ref{lem:unbiased} give
\[
\E\!\Bigl[\bigl\|\sum_{k}p_k\sum_{e=0}^{E-1}
      (g_{k}^{(t,e)}-\nabla f_k(x_{k}^{(t,e)}))\bigr\|^{2}\!\Big|\mathcal{F}_t\Bigr]
\le \sum_{k}p_k^{2}\,E\,\sigma^{2}
\le \frac{E\sigma^{2}}{K},
\tag{9}
\]
where we used $\sum_k p_k^2\le 1/K$ (equal participation).  

For the deterministic part, applying the triangle inequality and $\|a+b\|^2\le2(\|a\|^2+\|b\|^2)$,
\begin{align}
\Bigl\|\sum_{k}p_k\sum_{e=0}^{E-1}
      (\nabla f_k(x_{k}^{(t,e)})+\mu(x_{k}^{(t,e)}-x^{(t)}))\Bigr\|^{2}
&\le 2\Bigl\|\sum_{k}p_k\sum_{e=0}^{E-1}\nabla f_k(x_{k}^{(t,e)})\Bigr\|^{2}
   +2\mu^{2}\Bigl\|\sum_{k}p_k\sum_{e=0}^{E-1}(x_{k}^{(t,e)}-x^{(t)})\Bigr\|^{2} \nonumber\\
&\le 2E^{2}\Bigl(\sum_{k}p_k\|\nabla f_k(x_{k}^{(t,e)})\|\Bigr)^{2}
   +2\mu^{2}E^{2}\sum_{k}p_k\|x_{k}^{(t,e)}-x^{(t)}\|^{2}. \tag{10}
\end{align}
Using $L$‑smoothness and the fact that $f_k$ is $L$‑Lipschitz in its gradient,
$\|\nabla f_k(x_{k}^{(t,e)})\|\le L\|x_{k}^{(t,e)}-x^{(t)}\|+\|\nabla f_k(x^{(t)})\|$.
Dropping the non‑negative term $\|\nabla f_k(x^{(t)})\|$ yields an upper bound
proportional to $L^{2}\|x_{k}^{(t,e)}-x^{(t)}\|^{2}$. Consequently,
\[
\Bigl\|\sum_{k}p_k\sum_{e=0}^{E-1}
      (\nabla f_k(x_{k}^{(t,e)})+\mu(x_{k}^{(t,e)}-x^{(t)}))\Bigr\|^{2}
\le 2E^{2}(L^{2}+\mu^{2})\sum_{k}p_k\sum_{e=0}^{E-1}\|x_{k}^{(t,e)}-x^{(t)}\|^{2}.
\tag{11}
\]

Combining (9)–(11) into (8) gives
\[
\text{Term B}\;\le\;
\frac{2E\sigma^{2}}{K}
\;+\;4E^{2}(L^{2}+\mu^{2})\sum_{k}p_k\sum_{e=0}^{E-1}\|x_{k}^{(t,e)}-x^{(t)}\|^{2}.
\tag{12}
\]

\medskip
\noindent\textbf{[Step 6] Relate the drift $\|x_{k}^{(t,e)}-x^{(t)}\|$ to the stepsize.}  
From (\ref{eq:local_update}) and the bound $\eta_t\le 1/(L+\mu)$,
\begin{align}
\|x_{k}^{(t,e+1)}-x^{(t)}\|
&= \bigl\|x_{k}^{(t,e)}-x^{(t)} -\eta_t(g_{k}^{(t,e)}+\mu(x_{k}^{(t,e)}-x^{(t)}))\bigr\| \nonumber\\
&\le (1-\mu\eta_t)\|x_{k}^{(t,e)}-x^{(t)}\| + \eta_t\|g_{k}^{(t,e)}\| .
\end{align}
Taking expectation and using $\E\|g_{k}^{(t,e)}\|^{2}\le 2L^{2}\|x_{k}^{(t,e)}-x^{(t)}\|^{2}+2\sigma^{2}$ (by smoothness and variance bound) yields
\[
\E\bigl[\|x_{k}^{(t,e+1)}-x^{(t)}\|^{2}\mid\mathcal{F}_t\bigr]
\le (1-\mu\eta_t)^{2}\|x_{k}^{(t,e)}-x^{(t)}\|^{2}
      +\eta_t^{2}\bigl(2L^{2}\|x_{k}^{(t,e)}-x^{(t)}\|^{2}+2\sigma^{2}\bigr).
\]
Since $\eta_t\le 1/(L+\mu)$, the coefficient of $\|x_{k}^{(t,e)}-x^{(t)}\|^{2}$ is at most
$1-\frac{\mu}{L+\mu}\le 1$. Unrolling the recursion for $e=0,\dots,E-1$ gives
\[
\E\bigl[\|x_{k}^{(t,e)}-x^{(t)}\|^{2}\mid\mathcal{F}_t\bigr]
\le e\,\eta_t^{2}\sigma^{2},\qquad\forall e\ge1 .
\tag{13}
\]

Summing over $e$ and averaging over $k$:
\[
\sum_{k}p_k\sum_{e=0}^{E-1}\E\!\bigl[\|x_{k}^{(t,e)}-x^{(t)}\|^{2}\mid\mathcal{F}_t\bigr]
\le \eta_t^{2}\sigma^{2}\sum_{e=0}^{E-1}e
= \eta_t^{2}\sigma^{2}\frac{E(E-1)}{2}
\le \frac{\eta_t^{2}\sigma^{2}E^{2}}{2}.
\tag{14}
\]

\medskip
\noindent\textbf{[Step 7] Assemble the descent inequality.}  
Insert the lower bound (7) for Term A and the upper bound (12) for Term B
into (3), then take total expectation.
Using (14) to replace the drift sums yields
\begin{align}
\E\!\bigl[F(x^{(t+1)})\bigr]
&\le \E\!\bigl[F(x^{(t)})\bigr]
   -\eta_t\frac{E}{2}\E\!\bigl[\|\nabla F(x^{(t)})\|^{2}\bigr]
   +\eta_t\frac{L^{2}+\mu}{2}\underbrace{
          \E\!\bigl[\sum_{k}p_k\sum_{e=0}^{E-1}\|x_{k}^{(t,e)}-x^{(t)}\|^{2}\bigr]}_{\le \frac{\eta_t^{2}\sigma^{2}E^{2}}{2}}
   \nonumber\\
&\quad +\frac{L\eta_t^{2}}{2}
   \biggl(
        \frac{2E\sigma^{2}}{K}
        +4E^{2}(L^{2}+\mu^{2})\underbrace{
          \E\!\bigl[\sum_{k}p_k\sum_{e=0}^{E-1}\|x_{k}^{(t,e)}-x^{(t)}\|^{2}\bigr]}_{\le \frac{\eta_t^{2}\sigma^{2}E^{2}}{2}}
   \biggr) .
\end{align}
Simplifying and discarding higher‑order $\eta_t^{4}$ terms (which are
non‑negative) gives
\begin{equation}
\E\!\bigl[F(x^{(t+1)})\bigr]
\le \E\!\bigl[F(x^{(t)})\bigr]
   -\frac{E\eta_t}{2}\E\!\bigl[\|\nabla F(x^{(t)})\|^{2}\bigr]
   +\underbrace{\frac{L\sigma^{2}\eta_t^{2}E}{K}}_{\text{Variance term}}
   +\underbrace{L\mu^{2}\eta_t^{2}E^{2}}_{\text{Drift due to $\mu$}} .
\tag{15}
\end{equation}

\medskip
\noindent\textbf{[Step 8] Telescoping over $t=0,\dots,T-1$.}  
Sum (15) from $t=0$ to $T-1$ and rearrange:
\[
\frac{E}{2}\sum_{t=0}^{T-1}\eta_t\,\E\!\bigl[\|\nabla F(x^{(t)})\|^{2}\bigr]
\le F(x^{(0)})-F^\star
   +\sum_{t=0}^{T-1}\Bigl(
        \frac{L\sigma^{2}\eta_t^{2}E}{K}
        +L\mu^{2}\eta_t^{2}E^{2}
      \Bigr).
\]
Divide by $T$ and use $\eta_{\min}\le\eta_t\le\eta_{\max}$:
\[
\frac{E\eta_{\min}}{2}\Phi_T
\le \frac{F(x^{(0)})-F^\star}{T}
   +\frac{L\sigma^{2}\eta_{\max}^{2}E}{K}
   +L\mu^{2}\eta_{\max}^{2}E^{2}.
\]
Multiplying by $2/(E\eta_{\min})$ yields the claimed bound:
\[
\Phi_T\le
\frac{2\bigl(F(x^{(0)})-F^\star\bigr)}{\eta_{\min}T}
+\frac{2L\sigma^{2}\eta_{\max}E}{K}
+\frac{2\mu^{2}L\eta_{\max}E^{2}}{K}.
\tag{16}
\]

\medskip
\noindent\textbf{[Step 9] Choice of stepsize.}  
Set $\eta_t=\eta_0/\sqrt{t+1}$, then $\eta_{\min}\asymp 1/\sqrt{T}$,
$\eta_{\max}\asymp 1$. Plugging into (16) gives
\[
\Phi_T = O\!\Bigl(\frac{1}{\sqrt{T}}\Bigr) .
\]
Thus the algorithm attains the standard non‑convex sublinear rate,
with additional terms that vanish as $T\to\infty$ provided $E$ and $\mu$ are
chosen modestly.

\medskip
\noindent\textbf{[Step 10] Conclusion.}  
Inequality (16) is exactly the statement of Theorem~\ref{thm:main};
the proof is complete. \qed

%%%%%%%%%%%%%%%%%%%%%%%%%%%%%%%%%%%%%%%%%%%%%%%%%%%%%%%%%%%%%%%%%%%%%%%%%%%%%%
\section*{Rate Derivation (With Constants)}
Collecting the constants appearing in (16):
\[
C_{1}=2\bigl(F(x^{(0)})-F^\star\bigr),\qquad
C_{2}=2L\sigma^{2}E/K,\qquad
C_{3}=2L\mu^{2}E^{2}/K .
\]
Hence for any $T\ge1$,
\[
\boxed{\;
\Phi_T\le \frac{C_{1}}{\eta_{\min}T}+C_{2}\eta_{\max}+C_{3}\eta_{\max}
\;}
\]
which makes explicit how the heterogeneity penalty $\mu$ and the
communication interval $E$ influence the final rate.

%%%%%%%%%%%%%%%%%%%%%%%%%%%%%%%%%%%%%%%%%%%%%%%%%%%%%%%%%%%%%%%%%%%%%%%%%%%%%%
\section*{Special Cases}
\begin{enumerate}[label=(\alph*)]
\item \textbf{Convex $F$.} If each $f_k$ is convex, the same derivation yields a
      bound on the optimality gap $\E[F(x^{(T)})-F^\star]$ of order
      $O(1/\sqrt{T})$.
\item \textbf{Full‑batch local steps ($\sigma=0$).} The variance term disappears,
      and the rate improves to $O(1/T)$ provided $\eta_t$ is constant and
      $\eta_t\le 1/(L+\mu)$.
\item \textbf{No proximal term ($\mu=0$).} The bound reduces to the classical
      Local SGD result; the drift term becomes $O(E^{2})$ instead of $O(\mu^{2}E^{2})$.
\end{enumerate}

%%%%%%%%%%%%%%%%%%%%%%%%%%%%%%%%%%%%%%%%%%%%%%%%%%%%%%%%%%%%%%%%%%%%%%%%%%%%%%
\section*{Discussion}
The proof shows that the proximal regularizer in FedProx mitigates client drift
by converting the $E^{2}$ drift penalty into $ \mu^{2}E^{2}$; when data are
highly heterogeneous a modest $\mu$ (e.g. $0.1$–$1$) dramatically reduces the
constant $C_{3}$, leading to faster practical convergence while preserving the
theoretical $O(1/\sqrt{T})$ rate. The analysis relies only on smoothness and
bounded variance, thus holds for a broad class of non‑convex deep learning
objectives.

%%%%%%%%%%%%%%%%%%%%%%%%%%%%%%%%%%%%%%%%%%%%%%%%%%%%%%%%%%%%%%%%%%%%%%%%%%%%%%
\begin{thebibliography}{9}
\bibitem{li2020federated}
T.~Li, A.~K. Sahu, A.~Talwalkar, and V.~Smith.
\newblock Federated learning: Challenges, methods, and future directions.
\newblock {\em IEEE Signal Processing Magazine}, 37(3):50--60, 2020.
\end{thebibliography}

\end{document}